\subsection{The reduced type III cap for the primitive E7 source}

\begin{statement}{P3F.37}{Proved}
Let $S$ be the minimal regular surface of
\[
 y^2=x^3+tx+1,\qquad G=(0,1),\qquad
 M=\mathcal O_S(2G-O),\qquad Q=G.
\]
Choose any one of its three nodal fibres as the simple whole-fibre
zero of the linearization. For every $m=4l+3$, $l\ge0$, the
explicit root and affine norm below give a reduced type III cap
at the type III* fibre. Its marked meridian is
$m q+\gamma$, where $q$ is the positive base meridian and
$\gamma$ is the outward suspension loop.
The resulting construction is a proper holomorphic family over a
punctured $r$-disc, with scalar $\lambda=r^m$. Every fibre is
orientation-preservingly diffeomorphic to the standard $S^6$.
The construction and topology asserted here do not assign the
oriented parameter-return class.
\end{statement}
\begin{proof}
\par\medskip\noindent\textbf{The source and its actual root.}
The three finite discriminant roots are the distinct roots of
$4t^3+27$ and give $I_1$ fibres. At infinity put
$q=1/t$, $x_0=q^2x$, $y_0=q^3y$. The exact equation is
\[
 y_0^2=x_0^3+q^3x_0+q^6,\qquad G=(0,q^3).
\]
It has type III*. The source is rational, as witnessed on a
nonempty open set by $t=(y^2-x^3-1)/x$. The marked section data
of P3F.2 give $h(G)=1/2$, $G\cdot O=0$, and the nonidentity
component of the $E_7$ fibre. The height discriminant identity
then proves that $G$ is primitive and that the Mordell--Weil group
has no torsion. The section $2G$ has identity component and height
two, hence $2G\cdot O=0$ on this rational surface. Therefore
$M$ has multidegree zero on every component. Restricting the
actual Hom line to $O$ gives its degree
\[
 (2G)\cdot O+1-G\cdot O+(-G)\cdot O=1.
\]
Its vertical divisor has zero degree on each component, so it is
a whole-fibre pullback. Its section can consequently have its
sole simple zero at the selected $I_1$. Denote its line map by
$B_0$; it covers $t_{-G}$, and the contracting lift is
$B_\lambda=\lambda B_0$.

For $m=4l+3$ and $k=3l+2$ make the actual substitution
\[
 q=z^m,\qquad x_0=z^{2k}X,\qquad y_0=z^{3k}Y.
\]
The minimal equation and its marked section become
\[
 Y^2=X^3+zX+z^{6(l+1)},\qquad G'=(0,z^{3(l+1)}).       \tag{F37.1}
\]
Its only surface singularity over zero is the ordinary double
point with quadratic part $Y^2-zX$. One point resolution gives
a smooth source with a reduced III fibre. For example, at $m=3$
the chart $Y=Xv$, $z=Xw$ is
\[
 v^2=X+w+X^4w^6.
\]
The components $X=0$ and $w=0$ meet with tangency order two.
For general $l$ the same chart has the corresponding higher
monomial and the same smoothness and component configuration.
The exact deck is
\[
 D(X,Y,z)=(\zeta_m^{-2k}X,\zeta_m^{-3k}Y,\zeta_m z).   \tag{F37.2}
\]
It extends on the resolution and fixes both component labels.
At $m=3$ its chart formula is
$(X,v,w)\mapsto(\zeta_3^2X,\zeta_3v,\zeta_3^2w)$.
The section $G'$ meets the nonidentity component, while $2G'$
meets the identity component. This also follows from composing
with the good fourth root: its marked limit is the nonzero fixed
two-torsion point, as proved by the coordinate comparison in the
reduced-additive root calculation.

\par\medskip\noindent\textbf{The source norm and the actual line.}
Multiplication by the odd integer $m$ is a bijection on the
central smooth Neron group. On its additive identity component
its derivative is $m$, and on its component group $\mathbb Z/2$
it is invertible. The analytic inverse-function theorem gives the
unique local section $R$ with $mR=G'$ and the prescribed central
point. Uniqueness gives $D(R)=R$. The source translation extends
over the whole minimal regular surface: resolve its graph, or
apply the resolved addition and canonical-form comparison used
in P3K.22. A remaining exceptional curve would have to be a
vertical $(-1)$-curve on a relatively minimal source. Thus
\[
                   a_0=t_{-R}D,\qquad a_0^m=t_{-G'}.
\]
It exchanges the two III components.

Put $N=\mathcal O(2G'-O)$. Pull $M$ to the resolved fibre
product first. Its degree on a curve mapping to a point is zero,
and its degree on the transform of every source-fibre component
is zero as well. It therefore descends through each actual
$(-1)$-contraction of the relative minimalization, by the blowup
Picard decomposition. Its comparison with $N$ has a vertical
divisor of degree zero on the two III components. The matrix
\[
                    \begin{pmatrix}-2&2\\2&-2\end{pmatrix}
\]
has integral kernel $\mathbb Z(1,1)$, so this difference is an
integer whole fibre, principal on the disc. This identifies the
line, including its descended deck lift; that lift has $D^m=1$
and commutes with the pullback $B_0$.

An invertible line lift $L$ covering $t_{-R}$ can be constructed
without discarding its vertical divisor. For $P=2G'$, the chord
through $P,-R$, divided by the vertical function through $P-R$,
has fibre divisor
\[
                        P+(-R)-(P-R)-O.
\]
This is the divisor of $\operatorname{Hom}(t_R^*N,N)$ in the
specified direction. Its remaining vertical divisor has degree
zero on both components. The displayed matrix makes it an
integer whole fibre, removed by the corresponding base power.
The tangent formulas at coincident points give the same
meromorphic section and its extension.

Deck conjugation acts on the rank-one free module of such Hom
sections. Let $\chi$ be its central character. The finite average
\[
 L_{\rm eig}=\frac1m\sum_{j=0}^{m-1}
                 \chi^{-j}D^jLD^{-j}
\]
is nonzero at the central fibre and invertible after shrinking.
It satisfies $DL_{\rm eig}D^{-1}=\chi L_{\rm eig}$ and
$\chi^m=1$. Both $L_{\rm eig}^m$ and $B_0$ cover $t_{-G'}$.
Their ratio is a nonvanishing base function: scalar ratios are
constant on compact connected fibres. Deck invariance makes it
a function $u(q)$. Consequently
\[
 (L_{\rm eig}D)^m=\chi^{m(m-1)/2}u(q)B_0.
\]
Divide by an analytic $m$-th root of this base unit. The resulting
line map $A_0$ satisfies
\[
                   A_0^m=B_0,\qquad A_0B_0=B_0A_0.     \tag{F37.3}
\]
The character factor is one for this odd $m$, but was retained
before normalization. With $\lambda=r^m$ and $A_r=rA_0$ the
identity is exactly $A_r^m=B_\lambda$.

On a smaller compact source tube the norm ratios of $A_0$ and
its inverse are bounded. For small $|r|$, $A_r$ is strictly
contracting. Thus $N^\times/\langle A_r\rangle$ is a smooth
proper cap, with its compact norm shell as a fundamental region.
The line circle is free everywhere on this cap: equality of a
unit rotation and an iterate $A_r^j$ forces $j=0$ by strict
contraction, and then forces the rotation to be the identity.
This includes both components and their tangency point.

\par\medskip\noindent\textbf{The annular map and marked meridian.}
Use an outward increasing compatible log-norm coordinate $u$.
The actual cap identification is
\[
       (p,z,u)\sim(a_0p,\zeta_mz,u-1).
\]
The positive suspension generator is $h=A_r^{-1}$, and the
ordinary outward generator is $\gamma=B_\lambda^{-1}=h^m$.
Choose the positive root cut crossing to be the single map
$A_r$, without an additional $B_\lambda$ iterate. In real
angular coordinates its marking is
\[
 \frac{\arg z}{2\pi}=\frac{\arg q}{2\pi m},\qquad
 u_{\rm cap}=m u_{\rm ordinary}-\frac{\arg q}{2\pi}.
\]
Thus $q\mapsto-h$, $\gamma\mapsto mh$, and the killed
meridian is exactly $mq+\gamma$.

Here is a holomorphic realization of that marking. On a simply
connected sector a degree-zero source line has the constant-factor
presentation
\[
 (a,V)\sim(a+1,V),\qquad
 (a,V)\sim(a+\tau,e^{2\pi i\beta}V).
\]
The multiplication $(a,V)(a',V')=(a+a',VV')$ is its actual
commutative extension group. A line map covering a translation
lifts with a doubly periodic holomorphic multiplier, hence a
constant multiplier. The annular cap is a torsor for this same
group, with its cocycle relation precisely (F37.3). Lift the
cocycle to the holomorphic Lie bundle. Its integral period
two-cocycle is a coboundary because the annulus retracts onto a
circle. Remove it, and solve the resulting Lie cocycle using
$H^1=0$ for a holomorphic vector bundle on the Stein annulus.
Translation by the resulting zero-cochains gives mutually inverse
holomorphic gluing maps. The same exact sequence makes the
component map onto $H^1$ of the period local system surjective.
Its gamma class is primitive: the invariant integral quotient
$\psi$ takes gamma to one and the other periods to zero.
Choose the component with cut crossing $A_r$ itself; an additional
$B_\lambda^j$ is removed by the section of gamma winding $j$.
All source and line-circle phases remain in these translations.

Fix the arrow $F_c:\text{ordinary}\to\text{cap}$ and use the
physical character $V=e^{2\pi i\zeta_2}$, with
$c=\log\lambda/(2\pi i)$. Write $s=\log q/(2\pi i)$.
Its Lie translation obeys $g_c(s+1)=Dg_c(s)+a_c$.
Solving once at $c_0$ and adding $(c-c_0)s\delta/m$ gives its
exact parameter dependence. On a real angular cut its multiplier
is $r^s$, which is the norm marking displayed above. Since the
cap chart is identical at $c$ and $c+m$, the inverse closing map
is $F_c^{-1}F_{c+m}$, with multiplier $q$ near infinity.
Globally its multiplier is
\[
                    (t-t_c)^{-1}=\frac{q}{1-t_cq}.
\]
The positive parameter deck is its inverse. At the chosen cusp
multiplication by $t-t_c$ is, modulo $B_\lambda$, translation
by $+G$ and an invertible scalar unit. Both extend in the actual
nodal and Mumford charts, as does the inverse. At the other
fillings it is a unit; the displayed unit at infinity extends
over the III cap. The maps agree on their specified annuli.
They descend the construction to $r=e^{2\pi ic/m}$, and compact
norm shells give properness on each compact parameter annulus.

\par\medskip\noindent\textbf{The source-complement lattice.}
Use the integral divisor basis $O,F,C_0,\ldots,C_6,G$, where the
$E_7$ Cartan graph has edges
\[
 (0,1),(1,2),(2,3),(3,4),(4,5),(2,6),
\]
and $G$ meets $C_5$. The inverse Cartan column is
\[
 R_1=(1,2,3,5/2,2,3/2,3/2)^t.
\]
The full source Gram matrix has determinant $-1$, with
$O^2=G^2=-1$, $OF=GF=1$ and $GC_j=\delta_{j5}$.
These are therefore an integral basis. The full III* component
lattice, generated by $F$ and the seven $C_j$, is primitive.
For $V=S\setminus\operatorname{int}N_{III^*}$, excision and
the unimodular intersection pairing give
\[
 H_1(V)=H_3(V)=0,\qquad H_2(V)=\mathbb ZF\oplus\mathbb ZM,
 \quad M=P_{2G}-O=2G-2O-F+2R_1.                       \tag{F37.4}
\]
Orthogonality to $F$ and the root block forces the $G$ coefficient
to be even, proving this integral kernel description. The actual
source suspension $T=t_{-G}$ acts by
\[
                         TF=F,\qquad TM=M-F.           \tag{F37.5}
\]
For example the exact section formula
$P_{-G}=2O+2F-G-2R_1$ verifies the second equality directly.

The three nodal vanishing cycles generate the full elliptic
lattice. Rank one would make their monodromy product have trace
two. If their rank-two span had index $d>1$, every pairwise
determinant would be divisible by $d$. The trace of the product
of the three transvections is then $2$ modulo $d^2$: its quadratic
terms are determinant squares and its cubic term is their product.
The product is inverse to the III* monodromy and has trace zero.
Thus $d^2\mid2$, so $d=1$. Van Kampen with the section over
the disc base now gives $\pi_1(V)=1$.

The actual circle orbit exterior is the mapping torus of $T$ on
$V$. Its inclusion of the chosen cusp uses the relative nodal
orbit-product map of P3F.9, so no cusp piece is omitted. Denote
it by $W_V$. Wang's sequence and (F37.5) give
\[
 \pi_1(W_V)=\mathbb Z\gamma,\quad H_1(W_V)=\mathbb Z,
 \quad H_2(W_V)=\mathbb Z[M],\quad H_3(W_V)=\mathbb Z[F],
 \quad H_4(W_V)=0.
\]
The compact middle pairing is zero: the $H_3$ generator is the
trace of $F$, and $F\cdot M=0$ in a source slice.

\par\medskip\noindent\textbf{Every needed primitive gluing map.}
For the source boundary $L=\partial N_{III^*}$,
\[
 H_1(L)=\mathbb Zq\oplus\mathbb Z/2\,\eta,
 \qquad H_2(L)=\mathbb ZF.
\]
The inverse Cartan column $R_1$ modulo the root lattice is the
nonzero class $\eta$. Thus $Tq=q-\eta$ and $T\eta=\eta$.
An equivalent marked three-torus bundle for $E=\partial W_V$
has matrix
\[
 J=\begin{pmatrix}0&1&-1\\-1&0&0\\0&0&1\end{pmatrix}
\]
over $q$, in the elliptic and gamma basis. Both $J-I$ and
$\bigwedge^2J-I$ have Smith factors $(1,1,0)$. Their primitive
invariant vector and degree-two coinvariant functional are
$(-1,1,2)^t$ and $(2,-1,-1)$, respectively. Hence
\[
 H_1(E)=\mathbb Zq\oplus\mathbb Z\gamma,
 \quad H_2(E)=\mathbb Zz_1\oplus\mathbb Zz_2,
 \qquad F=2z_1.                                      \tag{F37.6}
\]
Poincare--Lefschetz duality and the zero compact middle pairing
make $H_2(E)\to H_2(W_V)$ onto with primitive kernel. Since $F$
maps to zero, that kernel is $\mathbb Zz_1$.

The signed second column has an actual relative-Wang representative.
Let $D_O$ be the punctured zero section, with boundary $q$, and
choose a boundary two-chain $C$ with $\partial C=2(Tq-q)$.
Then $2(TD_O-D_O)-C$ is a cycle. Changing $C$ adds a multiple
of $F$, zero in the target. Closing the section discs gives
\[
                  2(P_{-G}-O)=-M+3F-2R_1.
\]
The last two terms are in the removed fibre lattice, so the
second column is exactly $-M$.

The cap orbit $W_U$ is the mapping torus of $a_0$ on the III
tube. That tube retracts onto two spheres meeting at one point,
and $a_0$ exchanges them. Thus
\[
 \pi_1(W_U)=\mathbb Zh,\qquad H_2(W_U)=\mathbb ZC,
 \qquad H_3(W_U)=\mathbb Z(C_0+C_1).
\]
The incoming elliptic fibre specializes to $C_0+C_1$, giving
$F\mapsto2C$. Equation (F37.6) gives $z_1\mapsto C$.
Write $z_2\mapsto\kappa C$. The actual degree-two
Mayer--Vietoris matrix is
\[
 D_2=\begin{pmatrix}0&-1\\-1&-\kappa\end{pmatrix},
 \qquad \det D_2=-1.                                 \tag{F37.7}
\]
The determinant-one boundary operation $z_2\mapsto z_2-\kappa z_1$
makes its lower-right entry zero and leaves the exterior column
$-M$ unchanged. It does not alter the degree-one meridian.

The cap norm already gave $q\mapsto-h$, $\gamma\mapsto mh$;
the exterior sends $q$ to zero and gamma to its generator.
Consequently the full degree-one matrix is
\[
              D_1=\begin{pmatrix}0&1\\1&-m\end{pmatrix},
              \qquad\det D_1=-1.                     \tag{F37.8}
\]
We specify the higher maps as well. A primitive basis of $H_3(E)$
is the oriented source boundary $L$ followed by the source-fibre
gamma trace $F\gamma$. In the three-torus bundle description,
the former projects to the primitive invariant $F$ over $q$,
and the latter is the fundamental three-torus in a base slice.
The exterior row is $(0,1)$: $L$ bounds in the source complement
and $F\gamma$ is its Wang generator. For the cap, use the actual
projection to its norm-circle coordinate. On $L$ it has degree
$-1$, and on $F\gamma$ it has degree $m$, by the marked formula
$u=m\gamma-\arg q/(2\pi)$. At a regular norm value the corresponding
source cycle is the generic elliptic fibre. Its actual specialization
is $C_0+C_1$, the primitive invariant generator of the cap's source
$H_2$. Thus the cap row is exactly $(-1,m)$; there is no further
three-class from the source tube, which retracts onto its
two-dimensional central fibre. The full degree-three matrix is
\[
 D_3=\begin{pmatrix}0&1\\1&-m\end{pmatrix},
 \qquad \det D_3=-1.                                  \tag{F37.9}
\]
In degree four the map is $D_4:\mathbb Z\to0$, since both
pieces have $H_4=0$. In degree zero it is the incidence column
$D_0=(1,-1)^t$. These specify every nonzero gluing degree.

The elliptic boundary loops die in both cyclic fundamental groups.
These exact inclusions give $\pi_1(Y)=1$ by van Kampen.
Equations (F37.7)--(F37.9), $D_4$, and $D_0$ give all integral
sphere homology groups: only $H_0(Y)=H_5(Y)=\mathbb Z$ remains.
The smooth five-dimensional classification already used in
P3K.22 identifies this orbit with the standard $S^5$.

The only fixed set of the line circle is the single index-one
cusp sphere, with its actual $\mathcal O(-1)^{\oplus2}$ chart
and primitive Hopf meridian. The integral Gysin and fixed-tube
Mayer--Vietoris calculation of P3K.22 therefore gives
$\pi_1(X)=0$ and $H^*(X;\mathbb Z)=H^*(S^6;\mathbb Z)$.
The equality $\Theta_6=0$ identifies its oriented smooth type
with the standard sphere. The exact source, Smith and gluing
matrices are independently certified in
\path{code/part3_root/reduced_iii_audit.py} and its saved output.
The complete marked gluing degrees are also exported in
\path{code/part3/reduced_iii_cap_maps.py}. The physical-character
cochain and origin arithmetic are certified in
\path{code/part3/reduced_iii_return.py}. The separate actual
parameter-return comparison, with its signed Thom map, is P3K.24.
\end{proof}
